\section*{Tài liệu tham khảo}
\addcontentsline{toc}{section}{Tài liệu tham khảo}

\begin{thebibliography}{99}

\bibitem{tstcc2021}
Eldele, E., Ragab, M., Chen, Z., Wu, M., Kwoh, C. K., Li, X., \& Guan, C. (2021). 
Time-series representation learning via temporal and contextual contrasting. 
In \textit{Proceedings of the 30th International Joint Conference on Artificial Intelligence (IJCAI 2021)}, pp. 2352--2359.

\bibitem{thukral2025}
Thukral, M., Haresamudram, H., \& Ploetz, T. (2025). 
Cross-domain HAR: Few-shot transfer learning for human activity recognition. 
\textit{ACM Transactions on Intelligent Systems and Technology}, 16(1), pp. 1--35.

\bibitem{li2026}
Li, Y., et al. (2026). 
Bridging domain and instance gaps: A prototype contrastive framework for robust human activity recognition. 
\textit{Neurocomputing}, 132868.

\bibitem{crosshar2024}
Tang, C., et al. (2024). 
CrossHAR: Generalizing Cross-dataset Human Activity Recognition via Hierarchical Self-Supervised Pretraining. 
\textit{IEEE Transactions on Knowledge and Data Engineering (TKDE)}, arXiv preprint arXiv:2310.14390.

\bibitem{swav2020}
Caron, M., Misra, I., Mairal, J., Goyal, P., Bojanowski, P., \& Joulin, A. (2020). 
Unsupervised learning of visual features by contrasting cluster assignments. 
\textit{Advances in Neural Information Processing Systems (NeurIPS 2020)}, 33, pp. 9912--9924.

\bibitem{pcl2021}
Li, J., Zhou, P., Xiong, C., \& Hoi, S. C. (2021). 
Prototypical contrastive learning of unsupervised representations. 
In \textit{Proceedings of the International Conference on Learning Representations (ICLR 2021)}.

\bibitem{simclr2020}
Chen, T., Kornblith, S., Norouzi, M., \& Hinton, G. (2020). 
A simple framework for contrastive learning of visual representations. 
In \textit{Proceedings of the International Conference on Machine Learning (ICML 2020)}, pp. 1597--1607.

\bibitem{vit2021}
Dosovitskiy, A., Beyer, L., Kolesnikov, A., Weissenborn, D., Zhai, X., Unterthiner, T., ... \& Houlsby, N. (2021). 
An image is worth 16x16 words: Transformers for image recognition at scale. 
In \textit{Proceedings of the International Conference on Learning Representations (ICLR 2021)}.

\bibitem{motionsense2018}
Malekzadeh, M., Clegg, R. G., Cavallaro, A., \& Haddadi, H. (2018). 
Protecting sensory data against sensitive inferences. 
In \textit{Proceedings of the 11th ACM International Conference on Web Search and Data Mining (WSDM 2018)}, pp. 402--410.

\bibitem{ucihar2013}
Anguita, D., Ghio, A., Oneto, L., Parra, X., \& Reyes-Ortiz, J. L. (2013). 
A public domain dataset for human activity recognition using smartphones. 
In \textit{European Symposium on Artificial Neural Networks, Computational Intelligence and Machine Learning (ESANN 2013)}.

\bibitem{hhar2015}
Stisen, A., Blunck, H., Bhattacharya, S., Prentow, T. S., Kjærgaard, M. B., Dey, A. K., ... \& Jensen, C. S. (2015). 
Smart devices are different: Assessing and mitigating mobile sensing heterogeneity for activity recognition. 
In \textit{Proceedings of the 13th ACM Conference on Embedded Networked Sensor Systems (SenSys 2015)}, pp. 127--140.

\end{thebibliography}
